%
% teil3.tex -- Beispiel-File für Teil 3
%
% (c) 2020 Prof Dr Andreas Müller, Hochschule Rapperswil
%
% !TEX root = ../../buch.tex
% !TEX encoding = UTF-8
%
\section{3. Keplersches Gesetz mittels komplexer Integration}
Nach all den erarbeiten Grundlagen können wir nun endlich beginnen das 3. Keplersche Gesetz mittels komplexer Integration herzuleiten.
Wir gehen dafür von der Gleichung --Verweis einfügen-- aus welche in Abschnit --Verweis einfügen-- hergeleitet wurde. 
\[
T
=
\frac{m \cdot p^2}{L} \int_{0}^{2\pi} \frac{1}{(1+E \cdot cos(\theta))^2}\,d\theta
\]
Von dieser Gleichung interessierieren wir uns von nun an nur für das Integral.
Dieses könnte auch mittels Weierstrass-Substitution elegant gelöst werden.
\subsection{Substitution}
Die Grenzen 0 bis 2$\pi$ sind Indikator für eine Kontur in der komplexen Ebene.
Um diesem Indikator nachzugehen können wir eine Substitution machen:
\[
z
=
e^{i\theta}
\longrightarrow
cos(\theta)
=
\frac{e^{i\theta}+e^{-i\theta}}{2}
=
\frac{z+\frac{1}{z}}{2}
.
\]
Diese Substitution $z=e^i\theta$ beschreibt dabei eine komplexe Zahl auf dem komplexen Einheitskreis abhängig von $\theta$.
Um diese wieder in unser Integral einzusetzen müssen wir noch nach $dz$ umstellen:
\[
dz
=
\frac{d}{d\theta}e^{i\theta}\cdot d\theta
=
i \cdot e^{i\theta}\cdot d\theta
\Leftrightarrow
d\theta
=
\frac{dz}{i \cdot e^{i\theta}}
.
\]
Wenn wir dies nun in unser Integral einsetzen erhalten wir
\[
\int_{0}^{2\pi} \frac{1}{(1+E \cdot cos(\theta))^2}\,d\theta
=
\oint_C \frac{1}{(1+E \cdot \frac{z+\frac{1}{z}}{2})^2} \cdot \frac{1}{iz}\,dz
=
\frac{1}{i} \oint_C \frac{1}{z(1+E \cdot \frac{z+\frac{1}{z}}{2})^2}\,dz.
\]
Wobei $C = \left\lvert z\right\rvert = 1$ da wir einmal um den komplexen Einheitskreis herum integrieren.
\subsection{Vorbereitung für Cauchy-Integralsatz}
Unser Ziel ist es das Integral mit dem Cauchy-Integralsatz zu lösen.
Dieser hat gemäss --Satz 4.17-- folgende Form:
\[
f(z_0)
=
\frac{1}{2\pi i}
\oint_{\gamma} \frac{f(z)}{z-z_0}\,dz.
\]
Um den Integralsatz nutzen zu können muss das Innere unseres Integrals noch umgeformt werden um die Polstellen zu bestimmen:
\[
\frac{1}{i}\frac{1}{z(1+E \cdot \frac{z+\frac{1}{z}}{2})^2}
=
\frac{1}{i}\frac{z}{z^2(1+E \cdot \frac{z+\frac{1}{z}}{2})^2}
=
\frac{1}{i}\frac{z}{(z+\frac{Ez^2}{2}+\frac{E}{2})^2}
=
\]
\[
\frac{4}{i}\frac{z}{4(z+\frac{Ez^2}{2}+\frac{E}{2})^2}
=
\frac{4}{i}\frac{z}{(z+Ez^2+E)^2}
=
\frac{4}{iE^2}\frac{z}{(z^2+\frac{2}{E}z+1)^2}
\]
Mithilfe der ABC Formel kann das quadratische Polynom von $z$ noch faktorisiert werden:
\[
z^2+\frac{2}{E}z+1
=
0
\Leftrightarrow
z_0
=
\frac{-1}{E}\pm \sqrt{\frac{1}{E^2}-1}.
\]
Dies übernehmen wir in unser Integral:
\[
\frac{1}{i} \oint_C \frac{1}{z(1+E \cdot \frac{z+\frac{1}{z}}{2})^2}\,dz
=
\frac{4}{iE^2} \oint_C \frac{z}{(z+\frac{1}{E}+\sqrt{\frac{1}{E^2}-1})^2(z+\frac{1}{E}-\sqrt{\frac{1}{E^2}-1})^2}\,dz.
% Verweis möglich machen
\]
Nun haben wir die Polstellen gefunden und unser Integral erfolgreich umgeformt.
Jetzt muss verifiziert werden ob die Polstellen innerhalb unserer Kontur sind.
Dies führen wir mittels Grenzwert-Betrachtungen für $0 \leq E < 1$ durch.
\[
\lim_{E \to 1^-} \frac{-1}{E}+\sqrt{\frac{1}{E^2}-1}
=
\lim_{E \to 1^-} \frac{-1}{1^-}+\sqrt{\frac{1}{1^-}-1}
=
-1
\]
\[
\lim_{E \to 0^+} \frac{-1}{E}+\sqrt{\frac{1}{E^2}-1}
=
\lim_{E \to 0^+} \frac{-1}{0^+}+\sqrt{\frac{1}{0^+}-1}
=
0
\]
\[
\lim_{E \to 1^-} \frac{-1}{E}-\sqrt{\frac{1}{E^2}-1}
=
\lim_{E \to 1^-} \frac{-1}{1^-}-\sqrt{\frac{1}{1^-}-1}
=
-1
\]
\[
\lim_{E \to 0^+} \frac{-1}{E}-\sqrt{\frac{1}{E^2}-1}
=
\lim_{E \to 0^+} \frac{-1}{0^+}-\sqrt{\frac{1}{0^+}-1}
=
-\infty
\]
Die erste Nullstelle liegt innerhalb der Kontur, die zweite jedoch nicht.
Somit erhalten wir:
\[
\frac{4}{iE^2}
\oint_{\gamma} \frac{f(z)}{(z-z_0)^2}\,dz
,
\]
mit
\[
f(z)
=
\oint_C \frac{z}{(z+\frac{1}{E}+\sqrt{\frac{1}{E^2}-1})^2}\,dz
\qquad
\text{und $z_0=\frac{-1}{E}+\sqrt{\frac{1}{E^2}-1}$.}
\]
\subsection{Cauchy-Formel für Ableitungen}
In unserem Integral haben wir einen Pol zweiter Ordnung.
Auf Anblick sieht das nach einem Problem aus, da im Integralsatz keine Potenz im Nenner vorhanden ist.
Dies ist aber infolge --Satz 4.20-- kein Problem da folgende Gleichung gilt:
\[
f^{(n)}(z_0)
=
\frac{n!}{2\pi i}
\oint_C \frac{f(z)}{(z-z_0)^{n+1}}\,dz
\]
Um unser Integral zu lösen müssen wir nur noch die Gleichung von oben umformen und einsetzen.
Für $n=1$ gilt:
\[
f'(z_0)
=
\frac{1}{2\pi i}
\oint_C \frac{f(z)}{(z-z_0)^2}\,dz
\Leftrightarrow
\oint_C \frac{f(z)}{(z-z_0)^2}\,dz
=
2\pi i
\cdot
f'(z_0)
% Verweis möglich machen
\]
Als nächstes müssen wir die Funktion $f(z)$ mittels Quotientenregel ableiten.
\[
f'(z)
=
\frac{d}{dz}
\frac{z}{(z+\frac{1}{E}+\sqrt{\frac{1}{E^2}-1})^2}
=
\frac{1\cdot(z+\frac{1}{E}+\sqrt{\frac{1}{E^2}-1})^2-2z\cdot(z+\frac{1}{E}+\sqrt{\frac{1}{E^2}-1})}
{(z+\frac{1}{E}+\sqrt{\frac{1}{E^2}-1})^4}
=
\]
\[
\frac{z+\frac{1}{E}+\sqrt{\frac{1}{E^2}-1-2z}}{(z+\frac{1}{E}+\sqrt{\frac{1}{E^2}-1})^3}
=
\frac{\frac{1}{E}+\sqrt{\frac{1}{E^2}-1}-z}{(z+\frac{1}{E}+\sqrt{\frac{1}{E^2}-1})^3}
\]
Und noch $z_0$ in $f'(z)$ einsetzen:
\[
f'(z_0)
=
\frac{\frac{1}{E}+\sqrt{\frac{1}{E^2}-1}-z_0}{(z_0+\frac{1}{E}+\sqrt{\frac{1}{E^2}-1})^3}
=
\frac{\frac{1}{E}+\sqrt{\frac{1}{E^2}-1}+\frac{1}{E}-\sqrt{\frac{1}{E^2}-1}}
{(\frac{-1}{E}+\sqrt{\frac{1}{E^2}-1}+\frac{1}{E}+\sqrt{\frac{1}{E^2}-1})^3}
=
\]
\[
\frac{2}{E} \cdot \frac{1}{(2 \cdot \sqrt{\frac{1}{E^2}-1})^3}
=
\frac{1}{4E} \cdot (\frac{1}{\sqrt{\frac{1}{E^2}-1}})^\frac{3}{2}
=
\frac{1}{4E} \cdot (\frac{E^2}{1-E^2})^\frac{3}{2}
=
\frac{E^2}{4(1-E^2)^\frac{3}{2}}.
\]
Somit ist die rechte Seite der Gleichung --Verweis-- bestimmt.
\[
\oint_C \frac{f(z)}{(z-z_0)^2}\,dz
=
2\pi i
\cdot
f'(z_0)
=
2\pi i
\cdot
\frac{E^2}{4(1-E^2)^\frac{3}{2}}
\]
Wenn wir jetzt noch den Vorfaktor $\frac{4}{iE^2}$ von Gleichung --Verweis-- wieder einfügen vereinfacht es sich folgendermassen:
\[
\frac{4}{iE^2}
\cdot
\frac{2\pi i \cdot E^2}{4(1-E^2)^\frac{3}{2}}
=
\frac{2\pi}{(1-E^2)^\frac{3}{2}}
=
\int_{0}^{2\pi} \frac{1}{(1+E \cdot cos(\theta))^2}\,d\theta.
\]
Jetzt haben wir das Integral mittels komplexer Integration gelöst.
\subsection{Auflösen der Kepler-Gleichung}
Um das 3. Keplersche Gesetz fertig herzuleiten haben wir alle Bausteine zusammen.
Wir lösen die Gleichung --Verweis--:
\[
T
=
\frac{m \cdot p^2}{L} \int_{0}^{2\pi} \frac{1}{(1+E \cdot cos(\theta))^2}\,d\theta
\Leftrightarrow
T
=
\frac{m \cdot p^2}{L} \frac{2\pi}{(1-E^2)^\frac{3}{2}}.
\Leftrightarrow
T
=
\frac{m \cdot p^2}{m \sqrt{GMp}} \frac{2\pi}{(1-E^2)^\frac{3}{2}}
\]
\[
\Leftrightarrow
T
=
\frac{p^{\frac{3}{2}}}{\sqrt{GM}} \frac{2\pi}{(1-E^2)^\frac{3}{2}}
\Leftrightarrow
T^2
=
\frac{p^3}{GM} \frac{4\pi^2}{(1-E^2)^3}.
\]
Für den letzten Schritt muss noch folgendes beachtet werden: $p=\frac{b^2}{a}$ und $E=\sqrt{1-\frac{b^2}{a^2}}$.
\[
T^2
=
\frac{b^6}{a^3GM} \frac{4\pi^2}{(1-1+\frac{b^2}{a^2})^3}
\Leftrightarrow
T^2
=
\frac{b^6}{a^3GM} \frac{4\pi^2a^6}{b^6}
\Leftrightarrow
T^2
=
\frac{4\pi^2a^3}{GM}
\Leftrightarrow
\]
\[
\boxed{
\frac{T^2}{a^3}
=
\frac{4\pi^2}{GM}
=
const}
\]
Nun sehen wir das $T$ und $a$ die einzigen Parameter sind die verändert werden können.